%-------------------------------------------------------------------------------
\section{Analysis and Evaluation}
\label{sec:eval}
%-------------------------------------------------------------------------------

This section runs experiments with the methods and framework above and draws
quantitative conclusions. The internal pipeline metrics are reported in four
subsections in pipeline order: normative-rule extraction, lintability
determination, Lint Coverage Analysis, and code generation and semantic
equivalence. Beyond them, two mutually independent external validations
each check one half of the framework. \S\ref{subsec:lintability-ext} externally
validates normative-rule extraction and lintability determination against the
zlint maintainers' manual gold standard, and \S\ref{subsec:certdetect} externally
validates the generated lints by execution on real certificates.

\subsection{Experimental Setup}
\label{subsec:setup}

\begingroup\emergencystretch=2em\relax
All experiments use GPT-5.4 as the LLM for schema-constrained extraction and
coverage-alignment judging; the IR and code-accuracy audits are manual. The LLM
judges are configured at temperature~$0$. We target
RFC~5280 and CABF~BR because they are
the two most authoritative and foundational Web PKI normative sources for publicly
trusted TLS certificates. RFC~5280 defines the base X.509 certificate profile that
virtually every downstream normative rule inherits, and CABF~BR is the baseline that every
publicly trusted CA is mandated to meet; it also builds directly on RFC~5280, so
together they bound where compliance linting matters most. Normative-rule extraction was run on CABF~BR~v2.2.6. This was the latest
CABF~BR release when the experiments began.
\par\endgroup


The pipeline is evaluated in four stages reported below in execution order:
normative-rule extraction, lintability determination, Lint Coverage Analysis,
and code generation and semantic equivalence. The coverage analysis determines
which lintable rules are already implemented by zlint; only its uncovered
remainder enters the downstream code-generation experiment.

\subsection{Normative-Rule Extraction}
\label{subsec:rule-extraction}

The 1,804 raw capitalized RFC~2119 occurrences expand to 2,080 Layer-1
candidate records because nested clauses can inherit an obligation without
repeating its keyword.

The first-layer quality signal is G1 conservation; see Appendix~\ref{sec:saiv}. For the extracted candidates, the first conservation
identity is:
\[
\begin{array}{c @{\;} c @{\;} c @{\;} c @{\;} c @{\;} c @{\;} c}
  \underset{|\Rkw|}{2080} & = & \underset{|\RN|\ \text{noise}}{513} & + & \underset{|\RL|\ \text{lintable}}{278} & + & \underset{|\RU|\ \text{not-lintable}}{1289}
\end{array}
\]

\subsection{Lintability Determination}
\label{subsec:lintability-eval}

\begingroup\emergencystretch=2em\relax
The four-condition gate retains 278 single-certificate obligations and excludes
process, cross-certificate, runtime, external-semantic, MAY/OPTIONAL. These 278 lintable rules are the input to the zlint coverage
analysis below; no rule outside this set enters either coverage or code generation.
\par\endgroup


\begin{table}[!htbp]
\centering
\small
\setlength{\tabcolsep}{3pt}
\caption{Normative-rule extraction and lintability counts.}
\label{tab:rule-recall}
\begin{tabularx}{\columnwidth}{@{}X r r r@{}}
\toprule
Metric & CABF & RFC~5280 & Total \\
\midrule
Keyword occurrences & 1262 & 542 & 1804 \\
Recalled candidate records & 1440 & 640 & 2080 \\
Lintable rules & 176 & 102 & 278 \\
\bottomrule
\end{tabularx}
\end{table}

\subsection{Lint Coverage Analysis}
\label{subsec:coverage}

This analysis determines how many of the 278 lintable normative rules are
already implemented by an existing tool. We restrict the analysis to zlint,
the most widely deployed certificate linter, integrated into major CA issuance
and CT-monitoring pipelines. The criterion is deliberately strict. A lintable
normative rule counts as covered only when some cognate, or provenance-matched,
zlint lint genuinely implements it, not merely mentions the same field. We
accordingly label each normative rule full, partial, or none by whether such a
lint implements it fully, partially, or not at all, and count only full as
covered.

The matching procedure, including Algorithm~\ref{alg:coverage}, is
detailed in Appendix~\ref{app:coverage}.

\begin{table}[!htbp]
\centering
\small
\setlength{\tabcolsep}{3pt}
\caption{zlint coverage of the 278 lintable normative rules, by source.}
\label{tab:coverage}
\begin{tabularx}{\columnwidth}{@{}X r r r@{}}
\toprule
Item & CABF & RFC~5280 & Total \\
\midrule
lintable normative rules & 176 & 102 & 278 \\
zlint cognate lints & 170 & 122 & 292 \\
full coverage & 70 & 66 & 136 \\
uncovered code-generation domain & 106 & 36 & 142 \\
\bottomrule
\end{tabularx}
\end{table}

The deterministic aggregation yields the source-level covered and uncovered
counts in Table~\ref{tab:coverage}; only a full verdict contributes to the 136
covered rules, and every partial or none verdict remains in the 142-rule
code-generation domain $\RLunc$. The 292 cognate lints are lint instances, not
a covered-rule count.

The second conservation layer is the covered/uncovered split of the lintable
set in Eq.~\eqref{eq:split}:
\[
\begin{array}{c @{\;} c @{\;} c @{\;} c @{\;} c @{\;} c @{\;} c}
  \underset{|\RL|}{278} & = & \underset{|\RLcov|\ \text{covered}}{136} & + & \underset{|\RLunc|\ \text{uncovered}}{142}
\end{array}
\]
Both identities hold term by term, so the G1 residual $\Lrec = 0$ and the top
two conservation layers are deterministically recomputable. The
first layer is the three-way recall split of Eq.~\eqref{eq:cons}.
The 278 lintable normative rules further split by normative source into 176 CABF
and 102 RFC~5280, and the 142 uncovered normative rules $\RLunc$ are fed to the
downstream code generator.

\subsection{Code Generation and Semantic Equivalence}
\label{subsec:codegen-results}

Of the 142 uncovered rules, 128 have deterministic compilable lints, four use
LLM fallback, and ten have no delivered lint. Only manually verified generated
lints are included in the Bugzilla findings.

Human reviewers compared the IR or generated code with the selected source
text and applicable context; no LLM judged correctness. The frozen IR audit
found 122/142 source-matching records (85.92\%); the separately frozen code
audit found 115/128 source--code matches (89.84\%). For designated partial
rules, a correct nonempty subset on the same field and in the same direction
is accepted; none rules require the complete selected obligation.

\subsection{External Validation of Normative-Rule Extraction and Lintability
Determination}
\label{subsec:lintability-ext}

Normative-rule extraction and lintability determination need an independent external
reference. For this we bring in the CABF BR mapping table published by zlint
maintainers~\cite{zlint-mapping-sheet}. This table is manually annotated for each BR
normative rule by zlint experts as to whether it can be expressed by a certificate lint,
and the annotators are unaware of our extractor, IR, and lintability determination
module. Each Yes/No label in the table corresponds exactly to our joint output
of extracting that normative rule and determining whether it is lintable, and
agreement indicates that our normative-rule extraction and lintability determination match
human experts. The table covers two versions, BR~1.4.8~\cite{cabf-br148-md} and
BR~2.0.2~\cite{cabf-br202-md}, compared version-for-version against our
CABF-Server-1.4.8 extraction of 422 normative rules and the CABF-Server-2.0.2 extraction of 1024
normative rules, both produced by the same two-layer pipeline. Table~\ref{tab:extval}
reports the version-for-version agreement.

\begin{table}[!htbp]
\centering
\scriptsize
\setlength{\tabcolsep}{3pt}
\caption{Version-for-version external validation on BR~1.4.8 / BR~2.0.2, with
Cohen's $\kappa$ and 95\% confidence intervals.}
\label{tab:extval}
\begin{tabular*}{\columnwidth}{@{\extracolsep{\fill}}lrrrrr@{}}
\toprule
Version & Rows; Yes/No & Match cov. & Agreement & $\kappa$ & 95\% CI \\
\midrule
BR~1.4.8 & 122; 75/47 & 86.9\% & 91.5\% & 0.823 & [0.71, 0.93] \\
BR~2.0.2 & 250; 17/233 & 86.8\% & 96.8\% & 0.616 & [0.35, 0.88] \\
\bottomrule
\end{tabular*}
\end{table}

A small number of sheet rows fall outside our scope, so match coverage is near
87\% on both versions. These rows are title-case or lowercase descriptors such as
``Required if'', ``Deprecated'', or lowercase ``shall'' that RFC~2119 processing
correctly skips, together with, on the BR~2.0.2 sheet, cross-references or chapter
anchors with no matching extracted clause. On BR~2.0.2 the matched positive class is
very small, just 12 of 217 matched pairs, so a naive all-No baseline
already attains 94.5\% agreement. The more meaningful figure is therefore
$\kappa = 0.616$, substantial agreement on the Landis--Koch
scale~\cite{landis1977measurement}. BR~1.4.8 has more balanced classes, with 69
of 106 matched pairs positive, and yields $\kappa = 0.823$,
almost-perfect agreement and a tighter estimate of system quality.

\subsection{End-to-End Comparison with a Raw-Text Free-Form Baseline}
\label{subsec:baseline}

We compare the complete pipeline with a raw-text, free-form LLM-to-code workflow
on CABF~BR~v2.2.6 \S7.1.2, which contributes 141 of the 278 lintable normative
rules and 67 of the 142 rules in $\RLunc$. The baseline receives the verbatim text
of one subsection, enumerates the obligations it finds, and writes one zlint lint
for each obligation it judges lintable against a single parsed certificate. It
shares our model at temperature~0, our render-and-build gate, and our judging
protocol, and receives zlint~v3's registration skeleton and zcrypto's certificate
type, so failures cannot be mere unfamiliarity with zlint. It receives none of
the pipeline's apparatus: IR and rule records, the typed vocabulary $\mathcal{V}$,
atomic templates $\mathcal{A}$, the DSL grammar, the predicate-to-template map,
the deterministic renderer, and metadata binding. Both arms run end to end three
times on the same domain under identical prompts; the pipeline's Layer-1
extraction is deterministic. For semantic equivalence, the pipeline supplies its
deterministic code-semantics summary and the baseline raw Go source to the same
independent manual reviewer.

\begingroup\emergencystretch=2em\relax
To score them we enumerate mechanically all 496 obligation-class RFC~2119 keyword
occurrences in \S7.1.2 and adjudicate each, blind to both arms' outputs, on
whether it states a normative obligation and whether that obligation constrains
certificate content. This yields 493 obligations, 460 on certificate content. The
pipeline misses four of them, all table rows that Layer-1's regex matching failed
to segment; completing that logic would give 100\% recall.
\par\endgroup


Most of what the baseline accepts as lintable never compiles
(Table~\ref{tab:baseline}), and the accepted set is itself unstable from run to
run, so its denominator is an output of the baseline rather than a stated domain.
The pipeline splits its domain identically every run, 63 rules through the
deterministic main path and the remaining four through LLM fallback, and all
compile. Among the lints that do compile, the baseline's are likelier to be
marked non-equivalent by the manual reviewer to the obligation they claim to enforce. The variance
column follows the same line: what the pipeline does by machine does not move,
and the one stage it still leaves to manual review moves less than the baseline's.

\begin{table}[!htbp]
\centering
\scriptsize
\setlength{\tabcolsep}{3pt}
\caption{End-to-end comparison of a raw-text, free-form LLM-to-code baseline and
the complete pipeline.}
\label{tab:baseline}
\begin{tabular*}{\columnwidth}{@{\extracolsep{\fill}}lcccc@{}}
\toprule
Metric & \multicolumn{2}{c}{\shortstack{Raw-text\\free-form}} &
\multicolumn{2}{c}{\shortstack{Complete\\pipeline}} \\
 & \shortstack{Num.\\/$n$} & \shortstack{Rate\\(Var.)} &
 \shortstack{Num.\\/$n$} & \shortstack{Rate\\(Var.)} \\
\midrule
Extraction recall & \shortstack{$434/460$\\$432/460$\\$431/460$} &
  \shortstack{$94.0\%$\\$(1.56)$} & \shortstack{$456/460$\\$\times 3$} &
  \shortstack{$99.1\%$\\$(0.00)$} \\
\shortstack[l]{Code-generation\\success} & \shortstack{$99/344$\\$90/333$\\$98/343$} &
  \shortstack{$28.1\%$\\$(16.22)$} & \shortstack{$67/67$\\$\times 3$} &
  \shortstack{$100.0\%$\\$(0.00)$} \\
\shortstack[l]{Semantic\\equivalence} & \shortstack{$34/99$\\$40/90$\\$42/98$} &
  \shortstack{$40.5\%$\\$(11.56)$} & \shortstack{$53/67$\\$49/67$\\$47/67$} &
  \shortstack{$74.1\%$\\$(6.22)$} \\
\bottomrule
\end{tabular*}
\end{table}

\subsection{Bugzilla Certificate Detection}
\label{subsec:certdetect}

We used the manually verified generated lints alongside native zlint checks
to analyze one million Bugzilla certificates. The reported findings retain
only results from manually verified lints.

All retained classes identify specific Error/Warning conditions missing
from, or only partly implemented in, the native zlint version used.
A certificate may trigger native lints for other conditions. Every retained
occurrence has \texttt{notBefore} on or after its requirement's or
recommendation's effective date.
The date is checked per obligation; \texttt{notBefore} is a timing proxy.
After merging aliases and date filtering, Table~\ref{tab:bugzilla-findings}
contains 10 classes: 5 Error and 5 Warning. These comprise 543 Error and
2,939 Warning occurrences across 2,766 distinct certificates, including 540
with at least one Error. 

Missing DNS/IP SAN entries prevent SAN-based server identity matching.
Malformed SAN/CN labels can impair interoperability, although CN-only defects
need not affect SAN-based verification. Missing reserved policy identifiers
can disrupt policy-dependent validation. Prohibited OCSP \texttt{caIssuers}
violates the responder profile; responder AIA adds redundant information.
Missing \texttt{keyUsage} omits explicit restrictions on key use; other
constraints still apply. Extra OV subject attributes require verification
and maintenance. Policy qualifiers increase certificate size with limited
routine client value. Distribution points lacking HTTP/LDAP may hinder CRL
retrieval without alternatives.

\begin{table*}[t]
\centering
\small
\setlength{\tabcolsep}{4pt}
\caption{Source-reviewed findings in the one-million-file Bugzilla certificate scan after effective-date filtering.}
\label{tab:bugzilla-findings}
\begin{tabularx}{\textwidth}{@{}c X L{0.25\textwidth} r@{}}
\toprule
Severity & Finding class & Section & Occurrences \\
\midrule
Error & Present SAN has neither \texttt{dNSName} nor \texttt{iPAddress} & BR~7.1.2.7.12 & 3 \\
Error & SAN DNS label starts or ends with a hyphen & BR~7.1.2.7.12 & 4 \\
Error & Subscriber lacks a required reserved certificate policy identifier & BR~7.1.2.7.9 & 531 \\
Error & OCSP responder AIA contains prohibited \texttt{caIssuers} & BR~7.1.2.8.3 & 2 \\
Error & CN DNS label starts or ends with a hyphen & BR~7.1.4.3 & 3 \\
Warning & Subscriber lacks recommended \texttt{keyUsage} & BR~7.1.2.7.6 & 3 \\
Warning & OV Subscriber contains another discouraged subject attribute & BR~7.1.2.7.4 & 710 \\
Warning & Non-CA OCSP responder contains discouraged AIA & BR~7.1.2.8.2--3 & 2 \\
Warning & CA, Subscriber, or OCSP responder contains discouraged policy qualifiers & BR~7.1.2.7.9, 7.1.2.8.8, 7.1.2.10.5 & 2,210 \\
Warning & A DistributionPointName lacks an HTTP or LDAP URI & RFC~5280 4.2.1.13 & 14 \\
\bottomrule
\end{tabularx}
\end{table*}

